\documentclass[reqno,11pt]{amsart}
\usepackage{math327-lecture}
\usepackage{needspace}
\lectureinfo{13}{Oct. 5, 2026}{The Counting Formula and the Tower Law}

% Based on handwritten Lecture 10, Sep. 24, 2025.
\begin{document}
\makelecturetitle

\section{The counting formula for homomorphisms}


\begin{question} How many elements does $A_n$ have?

For $S_5$, the cycle types, their parity, and the number of permutations
of each type are listed below. Here $1$ denotes a fixed point.
\begin{center}
\normalfont
\begin{tabular}{ccl r}
\hline
Cycle type & Example & Parity & Count \\
\hline
$1^5$       & $1$             & Even & $1$  \\
$2\,1^3$   & $(12)$          & Odd  & $10$ \\
$2^2\,1$   & $(12)(34)$      & Even & $15$ \\
$3\,1^2$   & $(123)$         & Even & $20$ \\
$3\,2$     & $(123)(45)$     & Odd  & $20$ \\
$4\,1$     & $(1234)$        & Odd  & $30$ \\
$5$         & $(12345)$       & Even & $24$ \\
\hline
\end{tabular}
\end{center}
Thus $S_5$ has $60$ even and $60$ odd permutations, so $|A_5|=60$.

\newcommand{\rr}{\rightarrow}

By the way, how do we know that $(123)$ is even/odd? Remember our homomorphism $\sgn: S_n \rr \{ \pm 1\}$. Suppose $\sgn(123)=-1$, i.e. it is odd, then $\sgn\left( (123)^3 \right) = (-1)^3 = -1$, but it must also equal $1$ as $(123)^3=1$.

\end{question}

Let $\varphi:G\to H$ be a group homomorphism, and recall
\begin{align*}
  K=\ker(\varphi)&=\{x\in G:\varphi(x)=e_H\}\leq G,\\
  \im(\varphi)&=\{h\in H:\text{there is }a\in G\text{ with }\varphi(a)=h\}
  \leq H.
\end{align*}
Picture $G$ partitioned into the left cosets $K, g_1K, g_2K,\ldots$ of the
kernel. Since $\varphi(K)=\{e_H\}$, the whole coset $gK$ is sent to the
single element $\varphi(g)$:
\[
  \varphi(gk)=\varphi(g)\varphi(k)=\varphi(g)\qquad(k\in K).
\]
So $\varphi$ collapses each coset of $K$ to one point of $\im(\varphi)$.

\begin{proposition}\label{prop:cosets-image}
The map
\[
  G/K\longrightarrow\im(\varphi),\qquad gK\longmapsto\varphi(g),
\]
is a well-defined bijection.
\end{proposition}
\begin{proof}
\emph{Well defined:} if $g_1K=g_2K$, then by what we proved earlier $g_1^{-1}g_2 \in K.$ Hence $\varphi(g_1^{-1}g_2)=e$, i.e. $\varphi(g_1)=\varphi(g_2)$.

\emph{Surjective:} every element of $\im(\varphi)$ is $\varphi(g)$ for
some $g\in G$, and this is the image of $gK$.

\emph{Injective:} could we have $g_1K\neq g_2K$ but
$\varphi(g_1)=\varphi(g_2)$? No: if $\varphi(g_1)=\varphi(g_2)$, then
\[
  \varphi(g_2^{-1}g_1)=\varphi(g_2)^{-1}\varphi(g_1)=e_H,
\]
so $g_2^{-1}g_1\in K$. By the lemma on equality of cosets, this means
$g_1K=g_2K$.
\end{proof}

\begin{theorem}[Counting formula for homomorphisms]\label{thm:hom-count}
If $G$ is finite and $\varphi:G\to H$ is a homomorphism, then
\begin{equation}
  |G|=|\ker(\varphi)|\cdot|\im(\varphi)|.
\end{equation}
In particular, $|\im(\varphi)|$ divides both $|G|$ and $|H|$ (the latter
when $H$ is finite).
\end{theorem}
\begin{proof}
By Proposition~\ref{prop:cosets-image}, $|\im(\varphi)|=|G/\ker(\varphi)|$.
The counting formula applied to $\ker(\varphi)\leq G$ gives
$|G|=|\ker(\varphi)|\cdot|G/\ker(\varphi)|=|\ker(\varphi)|\cdot|\im(\varphi)|$.
Finally, $\im(\varphi)\leq H$, so Lagrange's theorem shows that
$|\im(\varphi)|$ divides $|H|$.
\end{proof}

\begin{example}[The alternating group]
For $n\geq2$, the sign homomorphism $\sgn:S_n\to\{\pm1\}$ is surjective
(a transposition has sign $-1$) and its kernel is $A_n$. Hence
\[
  n!=|A_n|\cdot2,\qquad |A_n|=\frac{n!}{2}.
\]
\end{example}

\section{Index in a tower of subgroups}

Recall the counting formula $|G|=[G:H]\,|H|$, where the index
$[G:H]=|G/H|$ is the number of left cosets of $H$ in $G$. (When we write
$[G:H]$ we usually mean that it is finite.)

\begin{theorem}[Tower law]\label{thm:tower}
Let $K\leq H\leq G$ be a tower of subgroups. Then
\begin{equation}
  [G:K]=[G:H]\,[H:K].
\end{equation}
\end{theorem}
\begin{proof}
Assume $[G:H]=m$ and $[H:K]=n$ are finite. Then
\[
  G=g_1H\sqcup g_2H\sqcup\cdots\sqcup g_mH\quad(g_i\in G),
  \qquad
  H=h_1K\sqcup\cdots\sqcup h_nK\quad(h_j\in H).
\]
Substituting the second decomposition into the first,
\[
  G=g_1(h_1K\cup\cdots\cup h_nK)\cup\cdots\cup
    g_m(h_1K\cup\cdots\cup h_nK)
   =\bigcup_{\substack{1\leq i\leq m\\1\leq j\leq n}}g_ih_jK.
\]
It remains to show that the $mn$ cosets $g_ih_jK$ are distinct, and hence
(being cosets) disjoint. Suppose $g_ih_jK=g_{i'}h_{j'}K$. Since
$h_jK\subseteq H$, we have $g_ih_jK\subseteq g_iH$ and similarly
$g_{i'}h_{j'}K\subseteq g_{i'}H$. So $g_iH\cap g_{i'}H\neq\varnothing$,
which forces $i=i'$. Cancelling $g_i$ then gives $h_jK=h_{j'}K$, so
$j=j'$. Therefore $G$ is the disjoint union of exactly $mn$ left cosets
of $K$, and $[G:K]=mn$.
\end{proof}

When $G$ is finite, the tower law also follows directly from the counting
formula:
\[
  [G:K]=\frac{|G|}{|K|}=\frac{|G|}{|H|}\cdot\frac{|H|}{|K|}=[G:H]\,[H:K].
\]
The proof above has the advantage that it works for infinite groups, as
long as the indices are finite.

\begin{example}
For $8\Z\leq2\Z\leq\Z$ we have $[\Z:2\Z]=2$ and $[2\Z:8\Z]=4$ (the cosets
are $8\Z$, $2+8\Z$, $4+8\Z$, $6+8\Z$), so
$[\Z:8\Z]=2\cdot4=8$, matching Lecture~11.
In $S_3$, the tower $\{1\}\leq\{1,x,x^2\}\leq S_3$ gives
$6=[S_3:\{1\}]=2\cdot3$.
\end{example}

%\section{Example: units modulo \texorpdfstring{$n$}{n} and Fermat's little theorem}
%
%Let $n\geq1$. The set $\Z/n\Z$ of residue classes modulo $n$ is a group
%under addition, with identity $\overline0$. Recall that
%\[
%  a\equiv b\pmod n\quad\text{means}\quad n\mid(b-a).
%\]
%Now consider the operation of \emph{multiplication} on this set,
%$\overline a\cdot\overline b=\overline{ab}$. It is not a group under
%multiplication (for instance $\overline0$ has no inverse), but we can
%ask: for which $a$ does the congruence
%\begin{equation}\label{eq:unit}
%  ax\equiv1\pmod n
%\end{equation}
%have a solution $x\in\Z/n\Z$?
%
%A solution means $ax-1=bn$ for some $b\in\Z$, i.e.
%\[
%  ax+(-b)n=1.
%\]
%Such an equation has an integer solution if and only if
%$\gcd(a,n)=1$: if $d$ divides $a$ and $n$ then it divides $1$;
%conversely, $\Z a+\Z n=\Z\gcd(a,n)$, so if $\gcd(a,n)=1$ then $1$ is an
%integer combination of $a$ and $n$. Note that $\gcd(a,n)$ depends only on
%$\overline a$, since $\gcd(a+kn,n)=\gcd(a,n)$.
%
%\begin{define}[Units modulo $n$]
%The \emph{group of units modulo $n$} is
%\[
%  (\Z/n\Z)^\times
%  =\{\overline a\in\Z/n\Z:\ ax\equiv1\ (\mathrm{mod}\ n)\text{ has a
%  solution}\}
%  =\{\overline a\in\Z/n\Z:\gcd(a,n)=1\},
%\]
%with the operation of multiplication.
%\end{define}
%
%This is not just a set but a group. The key point is closure: if
%$\overline a,\overline b\in(\Z/n\Z)^\times$, choose solutions
%$ax_1\equiv1$ and $bx_2\equiv1\pmod n$. Then $y=x_2x_1$ satisfies
%\[
%  (ab)y=a(bx_2)x_1\equiv ax_1\equiv1\pmod n,
%\]
%so $\overline{ab}\in(\Z/n\Z)^\times$. The remaining axioms are checked
%similarly.
%
%\begin{exercise}
%Complete the verification that $(\Z/n\Z)^\times$ is an abelian group:
%check associativity, that $\overline1$ is the identity, and that the
%solution $\overline x$ of \eqref{eq:unit} is an inverse of $\overline a$
%lying in $(\Z/n\Z)^\times$.
%\end{exercise}
%
%The order of this group is
%\[
%  |(\Z/n\Z)^\times|=\varphi(n)=\#\{1\leq a\leq n:\gcd(a,n)=1\},
%\]
%\emph{Euler's $\varphi$-function}. For example,
%$(\Z/8\Z)^\times=\{\overline1,\overline3,\overline5,\overline7\}$ and
%$\varphi(8)=4$.
%
%If $n=p$ is prime, every $a$ with $1\leq a\leq p-1$ is coprime to $p$, so
%$\varphi(p)=p-1$ and $(\Z/p\Z)^\times$ has order $p-1$. Lagrange's theorem
%now gives a classical result from number theory.
%
%\begin{theorem}[Fermat's little theorem]
%Let $p$ be prime and $a\in\Z$ with $p\nmid a$. Then
%\[
%  a^{p-1}\equiv1\pmod p.
%\]
%\end{theorem}
%\begin{proof}
%Since $p\nmid a$, $\overline a\in(\Z/p\Z)^\times$, a group of order
%$p-1$. By Lagrange's theorem the order $m$ of $\overline a$ divides
%$p-1$, say $p-1=mk$. Then $\overline a^{\,p-1}=(\overline a^{\,m})^k
%=\overline1$.
%\end{proof}
%
%The same argument applied to $(\Z/n\Z)^\times$ gives \emph{Euler's
%theorem}: if $\gcd(a,n)=1$, then $a^{\varphi(n)}\equiv1\pmod n$.
%
%\begin{remark}
%The group $(\Z/p\Z)^\times$ is in fact cyclic. A generator is called a
%\emph{primitive root} modulo $p$. For example, modulo $7$ the powers of
%$3$ are
%\[
%  3,\ 2,\ 6,\ 4,\ 5,\ 1,
%\]
%so $\overline3$ generates $(\Z/7\Z)^\times$. We will not prove this fact
%in general.
%\end{remark}
%
%\section{The correspondence theorem}
%
%Let $\varphi:G\to\mathcal G$ be a \emph{surjective} homomorphism of groups,
%and let $K=\ker(\varphi)\trianglelefteq G$.
%
%\begin{theorem}[Correspondence theorem]\label{thm:correspondence}
%\leavevmode
%\begin{enumerate}[label=\textbf{(\arabic*)}]
%  \item There is a bijective correspondence
%  \[
%    \{\text{subgroups }H\leq G\text{ with }K\leq H\}
%    \quad\longleftrightarrow\quad
%    \{\text{subgroups }\mathcal H\leq\mathcal G\},
%  \]
%  given by
%  \[
%    H\longmapsto\varphi(H)=\mathcal H,
%    \qquad
%    \mathcal H\longmapsto\varphi^{-1}(\mathcal H).
%  \]
%  \item Under this correspondence,
%  $H\trianglelefteq G$ if and only if
%  $\mathcal H\trianglelefteq\mathcal G$.
%  \item If $H$ is finite, then $|H|=|\mathcal H|\cdot|K|$.
%\end{enumerate}
%\end{theorem}
%
%Here $\varphi^{-1}(\mathcal H)=\{g\in G:\varphi(g)\in\mathcal H\}$ is the
%preimage. Part~(3) follows from what we already know: the restriction
%$\varphi|_H:H\to\mathcal H$ is a surjective homomorphism with kernel
%$H\cap K=K$, so the counting formula for homomorphisms (Theorem~\ref{thm:hom-count})
%gives $|H|=|K|\cdot|\mathcal H|$.

\end{document}
